\section{Module \texttt{date.c}}

Le module \texttt{date.c} encapsule le widget calendrier de GTK4. Il prend en charge la date initiale, l'affichage des noms de jours, les numéros de semaine et la remontée du signal de sélection.

\subsection{Fonction \texttt{create\_calendar}}

\begin{lstlisting}[language=C]
GtkWidget *create_calendar(const calendar_config *config) {
    GtkWidget *calendar = gtk_calendar_new();

    if (config == NULL) {
        return calendar;
    }

    if (config->selected_date != NULL) {
#if GTK_CHECK_VERSION(4, 20, 0)
        gtk_calendar_set_date(GTK_CALENDAR(calendar), config->selected_date);
#else
        gtk_calendar_select_day(GTK_CALENDAR(calendar), config->selected_date);
#endif
    }
    gtk_calendar_set_show_day_names(GTK_CALENDAR(calendar), config->show_day_names);
    gtk_calendar_set_show_week_numbers(GTK_CALENDAR(calendar), config->show_week_numbers);
    gtk_calendar_set_show_heading(GTK_CALENDAR(calendar), !config->no_month_change);
    apply_widget_style(calendar, &config->style);

    if (config->on_day_selected != NULL) {
        g_signal_connect(calendar, "day-selected", G_CALLBACK(config->on_day_selected), config->user_data);
    }

    return calendar;
}
\end{lstlisting}

\subsection{APIs GTK utilisées}

\begin{center}
\begin{tabular}{|l|p{8cm}|}
\hline
\textbf{API} & \textbf{Rôle} \\
\hline
\texttt{gtk\_calendar\_new()} & Crée l'instance du calendrier. \\
\hline
\texttt{gtk\_calendar\_set\_date()} & Positionne la date courante dans les versions récentes de GTK. \\
\hline
\texttt{gtk\_calendar\_select\_day()} & Assure la compatibilité avec les versions plus anciennes. \\
\hline
\texttt{gtk\_calendar\_set\_show\_week\_numbers()} & Active l'affichage du numéro de semaine. \\
\hline
\texttt{gtk\_calendar\_set\_show\_heading()} & Contrôle l'en-tête et la navigation mensuelle. \\
\hline
\texttt{g\_signal\_connect()} & Raccorde le signal de sélection de jour au callback. \\
\hline
\end{tabular}
\end{center}

\newpage
